\section{Related Work}
\label{sec:related}

\paragraph{Spatial benchmarks for VLMs.}
Recent work has shown that spatial reasoning remains a persistent weakness for vision-language models, even when their general multimodal scores appear strong. Dedicated benchmarks such as SpatialVLM \cite{chen2024spatialvlm}, VSR \cite{liu2023vsr}, and What'sUp \cite{kamath2023whatsup} probe relative position, orientation, and spatial language understanding more directly than generic VQA datasets. However, these benchmarks still evaluate models primarily through question-level correctness. They tell us whether a response matches the label, but not whether a model's responses can be assembled into a single scene-level geometric belief.

\paragraph{Controlled synthetic benchmarks.}
Ordinary-Bench is closer in spirit to controlled synthetic benchmarks such as CLEVR and CLEVRER \cite{johnson2017clevr,yi2020clevrer}, which trade ecological realism for exact ground truth and experimental control. This design choice is important for our setting: if the goal is to separate visual grounding from guessing and to test whether many local judgments imply a coherent spatial structure, then exact geometry, controlled complexity, and repeatable perturbations are more valuable than raw photorealism. We adopt this logic for dense ordinal probing rather than for compositional QA or event reasoning alone.

\paragraph{Language priors, hallucination, and evaluation mismatch.}
Several recent studies argue that apparent multimodal competence can be inflated by benchmark artifacts, textual shortcuts, or evaluation protocols that reward plausible language rather than genuine perception. BLINK highlights the gap between seeing and perceiving in multimodal models \cite{fu2024blink}, MMStar questions whether current VLM evaluation pipelines actually measure the intended capabilities \cite{chen2024mmstar}, POPE focuses on object hallucination \cite{li2023pope}, and VLInd-Bench explicitly measures language-prior contamination in VLM behavior \cite{lee2025vlind}. Ordinary-Bench complements these efforts by moving from answer-level artifacts to scene-level structure: our question is not only whether a model guesses correctly without vision, but whether its answers remain coherent when interpreted as a spatial belief.

\paragraph{Selective answering and reliability.}
A second line of work emphasizes that reliable multimodal evaluation should account for uncertainty, abstention, and response consistency. Reliable VQA argues that abstaining is often preferable to confidently answering incorrectly \cite{whitehead2022reliablevqa}; Selectively Answering Visual Questions and the consistency-based selective VQA framework of Khan and Fu further show that uncertainty and response stability are essential to black-box assessment \cite{eisenschlos2024selectivevqa,khan2024consistency}. Our protocol builds on this perspective. In addition to forced answering, Ordinary-Bench is designed to support selective-answering variants and to score structural consistency through transitivity and reciprocity, not just direct correctness.

\paragraph{Ordinal embedding and geometric reconstruction.}
The reconstruction component of Ordinary-Bench is grounded in work on ordinal embedding and geometric realizability. QRR constraints are naturally connected to ordinal embedding and non-metric multidimensional scaling \cite{agarwal2007gnmds,ma2018robustordinal}, while the question of whether a set of pairwise relations admits a valid geometric realization relates to Euclidean distance matrix theory \cite{dokmanic2015edm}. TRR constraints are closer to bearing and angle rigidity, which characterize when directional or angular relations determine a configuration up to known symmetries \cite{zhao2019bearingrigidity,chen2021anglerigidity}. We do not present a new generic reconstruction algorithm; rather, we use these ideas to define a benchmark analysis tool that tests whether VLM responses imply any coherent scene at all.
